% main_probability_metropolis.tex
\documentclass[aspectratio=169]{beamer}

% --- Language & encoding ---
\usepackage[utf8]{inputenc}
\usepackage[T1]{fontenc}
\usepackage[english,french]{babel}
\usepackage{lmodern}
\usepackage{microtype}

% --- Math & tables ---
\usepackage{amsmath,amssymb}
\usepackage{booktabs}
\usepackage{graphicx}
\usepackage{changepage}

% --- Metropolis Theme ---
\usetheme[numbering=fraction,progressbar=frametitle]{metropolis}
\setbeamertemplate{frame numbering}[fraction]

% --- Colors (same palette as prior lectures) ---
\definecolor{BootcampBlueDark}{HTML}{1A3C68}
\definecolor{TextBlack}{HTML}{000000}
\definecolor{AccentGold}{HTML}{DAA520}
\definecolor{Crimson}{HTML}{990000}
\definecolor{MatrixGreen}{HTML}{228B22}
\definecolor{MatrixRed}{HTML}{B22222}

\setbeamercolor{block title example}{bg=BootcampBlueDark!90!black,fg=white}
\setbeamercolor{block body example}{bg=BootcampBlueDark!5!white,fg=black}

% === Thick frame title band ===
\setbeamercolor{frametitle}{bg=BootcampBlueDark,fg=white}
\setbeamerfont{frametitle}{series=\bfseries,size=\large}
\setbeamertemplate{frametitle}{
  \nointerlineskip
  \begin{beamercolorbox}[wd=\paperwidth,ht=3.2ex,dp=1.4ex,leftskip=1em,rightskip=1em]{frametitle}
    \usebeamerfont{frametitle}\insertframetitle
  \end{beamercolorbox}
}

% === Custom Definition Block ===
\newenvironment{definitionblock}[1]{%
  \par\vspace{0.4em}%
  \begin{beamercolorbox}[
      wd=\textwidth,
      sep=0.4em,
      leftskip=0.3em,
      rightskip=0.6em,
      rounded=false,
      shadow=false
    ]{block title example}%
    \raisebox{0.15em}{\usebeamerfont{block title example}#1}%
  \end{beamercolorbox}%
  \vspace{-0.4em}%
  \begin{beamercolorbox}[
      wd=\textwidth,
      sep=1em,
      leftskip=0.6em,
      rightskip=0.6em,
      rounded=false,
      shadow=false
    ]{block body example}%
    \usebeamerfont{block body example}%
}{%
  \end{beamercolorbox}%
  \par\vspace{0.6em}%
}

% === Global text formatting ===
\setbeamercolor{normal text}{fg=TextBlack,bg=white}
\setbeamercolor{title}{fg=TextBlack}
\setbeamercolor{structure}{fg=BootcampBlueDark}
\setbeamercolor{progress bar}{fg=BootcampBlueDark,bg=gray!30}

% --- Course Info ---
\title{Probability Refresher}
\subtitle{Probability Spaces and Random Variables}
\author{Nathan Sauldubois}
\institute{NYU Tandon School, Finance and Risk Engineering}
\date{September 2021} % or \today
% \titlegraphic{\includegraphics[height=1.2cm]{your_logo.pdf}}

\metroset{
  sectionpage=progressbar,
  subsectionpage=progressbar
}

% --- Section & subsection transitions (plain) ---
\AtBeginSection{
  \frame[plain]{\sectionpage}
}
\AtBeginSubsection{
  \frame[plain]{\subsectionpage}
}

\begin{document}

\maketitle

% ===== Outline =====
\begin{frame}{Outline}
  \tableofcontents[hideallsubsections]
\end{frame}

% ===========================
% SECTION 0 — Motivation
% ===========================
\section{Motivation \& Reading}

\begin{frame}{Outline}
  \tableofcontents[currentsection, hideothersubsections]
\end{frame}

\begin{frame}{What is Probability Theory?}
\begin{itemize}
  \item Mathematical framework to model \textbf{uncertainty}.
  \item Describes \textbf{random experiments}, \textbf{events}, and their \textbf{likelihood}.
  \item Basis for statistics, stochastic processes, machine learning, finance, etc.
\end{itemize}
\end{frame}

\begin{frame}{Further Reading}
\begin{itemize}
  \item A. DasGupta, \textit{Probability for Statistics and Machine Learning}, Springer, 2011.
  \item Any introductory text on measure-theoretic probability.
\end{itemize}
\end{frame}

% ======================================
% SECTION 1 — Set Theory \& Probability Spaces
% ======================================
\section{Set Theory \& Probability Spaces}

\begin{frame}{Outline}
  \tableofcontents[currentsection, hideothersubsections]
\end{frame}

\subsection{Basic Set Theory}

\begin{frame}{Events as Sets}
\begin{definitionblock}{Event}
In probability theory, an \emph{event} is a subset of a sample space.
\end{definitionblock}
\pause 
\begin{itemize}
  \item Sets are denoted by capital letters \(A,B,C,\dots\)
  \item Example: rolling a dice, \(\Omega=\{1,2,3,4,5,6\}\)
    \begin{itemize}
      \item \(A=\{2\}\)
      \item \(B=\{2,4,6\}=\{\text{even result}\}\)
      \item \(C=\{1,2,3\}=\{\text{result}\le 3\}\)
    \end{itemize}
\end{itemize}
\end{frame}

\begin{frame}{Set Operations}
\begin{itemize}
  \item Complement: \(A^{c} = \{\omega : \omega \notin A\}\)
  \item Union: \(A \cup B = \{\omega : \omega \in A \text{ or } \omega \in B\}\)
  \item Intersection: \(A \cap B = \{\omega : \omega \in A \text{ and } \omega \in B\}\)
  \item Difference: \(A \setminus B = A \cap B^{c}\)
\end{itemize}
\pause 
\medskip
Example (dice roll, \(\Omega=\{1,\dots,6\}\)):
\begin{itemize}
  \item \(A^c = \{1,3,4,5,6\}\)
  \item \(A \cup C = C\)
  \item \(B \cap C = \{2\}\)
  \item \(B \setminus C = \{4,6\}\)
\end{itemize}
\end{frame}

\begin{frame}{Products of Sets \& Sequences}
\begin{itemize}
  \item Product of sets: \(A \times B = \{(a,b): a\in A,\,b\in B\}\).
  \pause 
  \item Infinite product: 
    \[
      A^{\mathbb{N}} = \{(a_1,a_2,\dots) : a_n \in A \text{ for all } n\}.
    \]
\end{itemize}
\pause 
Examples:
\begin{itemize}
  \item Two dice: \(\Omega = \{1,\dots,6\}\times\{1,\dots,6\}\).
  \item Infinite sequence of dice throws: \(\Omega=\{1,\dots,6\}^{\mathbb{N}}\).
\end{itemize}
\end{frame}

\begin{frame}{Sequences of Events}

\begin{definitionblock}{Sequence of Events}
A \textbf{sequence of events} is a family
\[
(A_n)_{n\ge 1}, \quad A_n \subset \Omega,
\]
where each \(A_n\) is an event in the same probability space.
\end{definitionblock}
\pause 
\vspace{-5mm}

\medskip

Two important types:

\begin{itemize}
  \item \textbf{Increasing sequence}
  $
    A_1 \subset A_2 \subset A_3 \subset \cdots
  $
  The event “eventually occurs” is
  $
    \bigcup_{n\ge 1} A_n
  $.
\pause 
  \item \textbf{Decreasing sequence}
  $
    B_1 \supset B_2 \supset B_3 \supset \cdots
  $
  The event “persists forever” is
  $
    \bigcap_{n\ge 1} B_n
  $
\end{itemize}

\end{frame}

\begin{frame}{Infinite Unions and Intersections}

Let \((A_n)_{n\ge1}\) be a sequence of events.
\pause 
\medskip

\textbf{Infinite union}
\[
\bigcup_{n\ge1} A_n
\]
is the event that \emph{at least one} of the events \(A_n\) occurs:
\[
\omega \in \bigcup_{n\ge1} A_n
\quad \Longleftrightarrow \quad
\exists\, n \ge 1 \text{ such that } \omega \in A_n.
\]
\pause 
\medskip

\textbf{Infinite intersection}
\[
\bigcap_{n\ge1} A_n
\]
is the event that \emph{all} events \(A_n\) occur:
\[
\omega \in \bigcap_{n\ge1} A_n
\quad \Longleftrightarrow \quad
\forall\, n \ge 1,\ \omega \in A_n.
\vspace{-5mm}
\]

\medskip
\pause 
\textbf{Interpretation:}
\begin{itemize}
  \item infinite union = ``occurs at least once''
  \item infinite intersection = ``occurs forever''
\end{itemize}

\end{frame}

\begin{frame}{Examples: Sequences of Events}

\textbf{Example 1 — Increasing events}

Let
$
A_n=\{\text{maximum of first $n$ dice throws is 6}\}.
$
Then:
$
A_1 \subset A_2 \subset A_3 \subset \cdots
.$
\pause 
Interpretation:
\begin{itemize}
  \item as \(n\) increases, it becomes easier to have seen a 6
  \item the event
  $
  \bigcup_{n\ge 1} A_n
  $
  means “a 6 appears at least once eventually”.
\end{itemize}

\medskip

\textbf{Example 2 — Decreasing events}

Let
$
B_n=\{\text{all first $n$ dice results are even}\}.
$ Then:
$
B_1 \supset B_2 \supset B_3 \supset \cdots
$.
Interpretation: the condition becomes stricter as \(n\) increases.

\end{frame}

\begin{frame}{\(\liminf\) and \(\limsup\) of Events}

\begin{definitionblock}{Limit Inferior and Superior of Events}
For a sequence of events \((A_n)_{n\ge 1}\),
\[
\liminf_{n\to\infty} A_n
= \bigcup_{n=1}^{\infty} \bigcap_{k\ge n} A_k,
\qquad
\limsup_{n\to\infty} A_n
= \bigcap_{n=1}^{\infty} \bigcup_{k\ge n} A_k.
\]
\vspace{-6mm}
\end{definitionblock}
\vspace{-6mm}
\medskip
\pause 
\textbf{Interpretation:}
\begin{itemize}
  \item \(\liminf A_n\): event occurs \emph{from some point on} (eventually always)
  \item \(\limsup A_n\): event occurs \emph{infinitely often}
\end{itemize}


\end{frame}

\begin{frame}{Examples of \(\liminf\) and \(\limsup\)}

Let \(A_n=\{\text{“the $n$-th coin toss is Head”}\}\).
\pause 
\begin{itemize}
  \item \(\limsup A_n\): Event that Head occurs infinitely often
  $
  \limsup_{n\to\infty} A_n
  = \{\text{H appears infinitely many times}\}.
  $
\pause 
  \item \(\liminf A_n\):
  Event that from some time onward all outcomes are Head
  $
  \liminf_{n\to\infty} A_n
  = \{\text{from some $N$ onward, every toss is H}\}.
  $
\end{itemize}
\pause 
\medskip

\textbf{Key facts:}
\begin{itemize}
  \item \(\liminf A_n \subseteq \limsup A_n\)
  \item for a fair coin:
    \[
      \mathbb{P}(\limsup A_n)=1,
      \qquad
      \mathbb{P}(\liminf A_n)=0
    \]
\end{itemize}

\end{frame}

\subsection{Sample Space, Events, Probability}

\begin{frame}{Sample Space \& Events}
\begin{definitionblock}{Sample Space}
A \emph{sample space} \(\Omega\) is a finite or infinite set representing all possible outcomes of a random experiment.
\end{definitionblock}
\pause 
\begin{itemize}
  \item Any subset \(A \subseteq \Omega\) is an \emph{event}, including \(\Omega\) itself and \(\emptyset\).
\end{itemize}
\pause 
\medskip
Examples:
\begin{itemize}
  \item dice roll: \(\Omega = \{1,\dots,6\}\)
  \item Infinite coin tossing: 
  \[
    \Omega = \{H,T\}^{\mathbb{N}} = \{(\omega_1,\omega_2,\dots) : \omega_i \in \{H,T\}\}.
  \]
\end{itemize}
\end{frame}

\begin{frame}{Basic Probability Measure}
\begin{definitionblock}{Probability Measure}
Let \(\Omega\) be a sample space. A probability measure \(\mathbb{P}\) assigns a number in \([0,1]\) to each event such that:
\pause 
\begin{itemize}
  \item \(\mathbb{P}(\emptyset) = 0\), \(\mathbb{P}(\Omega) = 1\);
  \item (Countable additivity) For every sequence of pairwise disjoint events \(A_1,A_2,\dots\),
  \[
    \mathbb{P}\Big(\bigcup_{n\ge1} A_n\Big) = \sum_{n\ge1} \mathbb{P}(A_n).
  \]
\end{itemize}
\pause 
\end{definitionblock}
The pair \((\Omega,\mathbb{P})\) is called a \emph{probability space}.
\end{frame}

\begin{frame}{Examples of Probability Measures}

\textbf{Example 1 — Fair dice}
\vspace{-3mm}
\[
\Omega=\{1,2,3,4,5,6\}, \qquad
\mathbb{P}(\{k\})=\frac{1}{6}.
\]
\pause 
Event example:
$
A=\{2,4,6\} \quad \Rightarrow\quad
\mathbb{P}(A)=\frac{3}{6}=\frac{1}{2}.
$
\pause 

\medskip
\textbf{Example 2 — Biased coin}
\vspace{-3mm}
\[
\Omega=\{H,T\}, \qquad
\mathbb{P}(H)=p,\quad \mathbb{P}(T)=1-p.
\]
\pause 
Event example:
$
A=\{H\} \Rightarrow \mathbb{P}(A)=p.
$
\pause 

\medskip
\textbf{Example 3 — Default of a firm}
\vspace{-3mm}
\[
\Omega=\{\text{Default},\text{No Default}\}
\]

Let
$
\mathbb{P}(\text{Default})=0.02,\qquad
\mathbb{P}(\text{No Default})=0.98.
$

\end{frame}

\begin{frame}{Examples of Measures}
\begin{itemize}
  \item \textbf{Uniform measure} on finite \(\Omega\):
  \[
    \mu(A) = \frac{\text{card}(A)}{\text{card}(\Omega)}.
  \]
  \pause 
  \item \textbf{Dirac measure} at point \(a\), denoted \(\delta_a\):
  \[
    \delta_a(A) = 
      \begin{cases}
        1 & \text{if } a\in A,\\
        0 & \text{otherwise.}
      \end{cases}
  \]
  \pause 
  \item \textbf{Lebesgue measure} \(\lambda\) on \(\mathbb{R}\):
  \[
    \lambda([a,b]) = b-a.
  \]
  Not a probability measure (total mass is infinite).
  \pause 
\end{itemize}
\end{frame}

\begin{frame}{Basic Properties of \(\mathbb{P}\)}
Let \((\Omega,\mathbb{P})\) be a probability space.
\begin{itemize}
  \item If \(A \subseteq B\), then \(\mathbb{P}(A) \le \mathbb{P}(B)\).
  \pause 
  \item For any event \(A\),
    \[
      \mathbb{P}(A^c) = 1 - \mathbb{P}(A).
    \]\pause 
  \item For any events \(A,B\),
    \[
      \mathbb{P}(A \cup B) = \mathbb{P}(A) + \mathbb{P}(B) - \mathbb{P}(A\cap B)
      \le \mathbb{P}(A)+\mathbb{P}(B).
    \]\pause 
\end{itemize}
\medskip
\pause 
\vspace{-6mm}
\begin{center}
\begin{tikzpicture}[scale=0.7]

  % Universe
  \draw (-3,-2) rectangle (3,2);
  \node at (-2.7,1.7) {$\Omega$};

  % Sets A and B
  \draw (-1,0) circle (1.5);
  \draw (1,0) circle (1.5);

  % Labels
  \node at (-1.6,0.8) {$A$};
  \node at (1.6,0.8) {$B$};
  \node at (0,0) {$A\cap B$};

\end{tikzpicture}
\end{center}
\end{frame}

\begin{frame}{Union Bound \& Total Probability}
\begin{itemize}
  \item \textbf{Union bound.} For any sequence \((A_n)_{n\ge1}\),
  \[
    \mathbb{P}\Big(\bigcup_{n\ge1} A_n\Big) \le \sum_{n\ge1} \mathbb{P}(A_n).
  \]
  \pause 
  \item \textbf{Law of total probability.} Let \(A\) be an event and \((B_n)_{n\ge1}\) a sequence of disjoint sets with \(\bigcup_{n\ge1} B_n = \Omega\). Then
  \[
    \mathbb{P}(A) = \sum_{n\ge1} \mathbb{P}(A \cap B_n).
  \]
  \pause 
\end{itemize}
\end{frame}

\begin{frame}{Examples: Union Bound and Total Probability}

\textbf{Union bound}

Throw a dice once.

$
A=\{1\},\qquad B=\{2\},\qquad C=\{3\}.
$
\pause 
Then
$
\mathbb{P}(A\cup B\cup C)=\frac{3}{6}
\le \frac{1}{6}+\frac{1}{6}+\frac{1}{6}.
$

\end{frame}

\subsection{Independence \& Conditioning}

\begin{frame}{Independence of Events}
\begin{definitionblock}{Independence}
Two events \(A\) and \(B\) are independent if
\[
  \mathbb{P}(A\cap B) = \mathbb{P}(A)\mathbb{P}(B).
\]
\pause 
Events \(A_1,\dots,A_n\) are (mutually) independent if for every subset
\(\{i_1,\dots,i_k\}\subset\{1,\dots,n\}\),
\[
  \mathbb{P}(A_{i_1}\cap\cdots\cap A_{i_k})
  = \mathbb{P}(A_{i_1})\cdots\mathbb{P}(A_{i_k}).
\]
\end{definitionblock}
\end{frame}

\begin{frame}{Examples of Independent Events}

\textbf{Example 1 — Two coin tosses}

Sample space:
\[
\Omega=\{HH,HT,TH,TT\}
\]
\pause 
Events:
\[
A=\{\text{first toss is H}\}=\{HH,HT\}
\]
\[
B=\{\text{second toss is H}\}=\{HH,TH\}
\]
\pause 
Then
\[
\mathbb{P}(A)=\mathbb{P}(B)=\frac{1}{2},
\qquad
\mathbb{P}(A\cap B)=\mathbb{P}(\{HH\})=\frac{1}{4}
\]
\pause 
Thus
\[
\mathbb{P}(A\cap B)=\frac{1}{4}
=\frac{1}{2}\cdot\frac{1}{2}
=\mathbb{P}(A)\mathbb{P}(B)
\]

\textbf{Conclusion:} \(A\) and \(B\) are independent.

\end{frame}


\begin{frame}{Conditional Probability}
\begin{definitionblock}{Conditional Probability}
Let \(A,B\) be events with \(\mathbb{P}(B)>0\). The conditional probability
of \(A\) given \(B\) is
\[
  \mathbb{P}(A \mid B)
  = \frac{\mathbb{P}(A\cap B)}{\mathbb{P}(B)}.
\]
\pause
\end{definitionblock}
\begin{itemize}
  \item Interpreted as ``probability of \(A\) knowing that \(B\) occurs''.
\end{itemize}
\pause
\medskip
\textbf{Multiplicative formula.} For events \(A_1,\dots,A_n\),
\[
\mathbb{P}\Big(\bigcap_{k=1}^n A_k\Big)
= \mathbb{P}(A_1)\,
  \mathbb{P}(A_2\mid A_1)\,
  \mathbb{P}(A_3\mid A_1\cap A_2)\cdots
  \mathbb{P}(A_n\mid A_1\cap\cdots\cap A_{n-1}).
\]
\end{frame}

\begin{frame}{Examples of Conditional Probability}

\textbf{Example — dice roll}

Let
$
A=\{\text{result is even}\}=\{2,4,6\},
\,\,
B=\{\text{result} > 3\}=\{4,5,6\}.
$
\pause
Then
\[
\mathbb{P}(A)=\frac{3}{6}=\frac{1}{2},
\qquad
\mathbb{P}(B)=\frac{3}{6}=\frac{1}{2}.
\]
\pause
Intersection:
\[
A\cap B=\{4,6\},
\qquad
\mathbb{P}(A\cap B)=\frac{2}{6}=\frac{1}{3}.
\]
\pause
Thus
\[
\mathbb{P}(A\mid B)
=\frac{\mathbb{P}(A\cap B)}{\mathbb{P}(B)}
=\frac{\frac{1}{3}}{\frac{1}{2}}
=\frac{2}{3}.
\]
\pause
\textbf{Interpretation:}

Given that the result is \(>3\), the probability of being even is \(2/3\).

\end{frame}

\begin{frame}{Total Probability \& Bayes in Terms of Conditioning}
\begin{definitionblock}{Law of Total Probability}
Let \(A\) be an event and \((B_n)_{n\ge1}\) a partition of \(\Omega\) (which means that $ \cup B_n = \Omega $ and $B_n \cap B_k = \emptyset$ for $k \neq n$) with \(\mathbb{P}(B_n)>0\). Then
\[
  \mathbb{P}(A) = \sum_{n\ge1} \mathbb{P}(A\mid B_n)\,\mathbb{P}(B_n).
\]
\pause
\end{definitionblock}
\begin{definitionblock}{Bayes' Formula (finite case)}
For events \(A,B\) with \(\mathbb{P}(B)>0\),
\[
  \mathbb{P}(A\mid B)
  = \frac{\mathbb{P}(B\mid A)\,\mathbb{P}(A)}{\mathbb{P}(B)}.
\]
\end{definitionblock}
\end{frame}


\subsection{Monotone Sequences of Events}

\begin{frame}{Monotone Sequences of Events}
Let \((\Omega,\mathbb{P})\) be a probability space.
\begin{itemize}
  \item An \emph{increasing} sequence: \(A_1 \subseteq A_2 \subseteq A_3 \subseteq \cdots\).
  \item A \emph{decreasing} sequence: \(B_1 \supseteq B_2 \supseteq B_3 \supseteq \cdots\).
\end{itemize}
\pause
\begin{definitionblock}{Continuity of Probability}
\begin{itemize}
  \item If \((A_n)\) is increasing,
    \[
      \mathbb{P}\Big(\bigcup_{n\ge1} A_n\Big) = \lim_{n\to\infty} \mathbb{P}(A_n).
    \]
    \pause
  \item If \((B_n)\) is decreasing,
    \[
      \mathbb{P}\Big(\bigcap_{n\ge1} B_n\Big) = \lim_{n\to\infty} \mathbb{P}(B_n).
    \]
\end{itemize}
\end{definitionblock}
\end{frame}

\begin{frame}{Example of a Monotone Sequence}

\textbf{Example (coin tossing).}

Consider an infinite sequence of fair coin tosses and define
\[
A_n = \{\text{at least one Head appears in the first $n$ tosses}\}.
\]
\pause
Then:
\[
A_1 \subset A_2 \subset A_3 \subset \cdots
\]
\pause
\begin{itemize}
  \item the event becomes less restrictive as \(n\) increases
  \item \(\bigcup_{n\ge1} A_n\) means:
  \[
  \{\text{a Head appears at least once eventually}\}
  \]
\end{itemize}

Since the coin is fair,
\[
\mathbb{P}(A_n)=1-\left(\frac12\right)^n
\quad\Longrightarrow\quad
\mathbb{P}\Big(\bigcup_{n\ge1} A_n\Big)=1.
\]

\end{frame}

\begin{frame}{Borel--Cantelli Lemma}

Let \((A_n)_{n\ge1}\) be a sequence of events.
\pause
\begin{definitionblock}{Borel--Cantelli Lemma (I)}
If
$
\sum_{n=1}^\infty \mathbb{P}(A_n) < \infty,
$
\pause
then
\[
\mathbb{P}\big(A_n \text{ infinitely often}\big) = 0.
\]
where 
\[
A_n \text{ infinitely often}
\;=\;
\limsup_{n\to\infty} A_n
=
\bigcap_{n=1}^{\infty} \bigcup_{k\ge n} A_k.
\vspace{-5mm}
\]
\pause
\end{definitionblock}

\textbf{Interpretation:}

If events are rare enough (summable probabilities),  
then only finitely many of them occur almost surely.

\end{frame}

\begin{frame}{Illustration of Borel--Cantelli}

\textbf{Example.}

Let
\[
A_n = \{\text{the $n$-th coin toss is Head}\}.
\]
\pause
Then:
\[
\mathbb{P}(A_n)=\frac12,
\qquad
\sum_{n=1}^\infty \mathbb{P}(A_n)=\infty.
\]
\pause
\medskip

\textbf{Conclusion:}
\[
\mathbb{P}\big(A_n \text{ infinitely often}\big)=1.
\]

\medskip

\textbf{Meaning:}
\begin{itemize}
  \item Heads occur infinitely many times almost surely
  \item this does \emph{not} contradict randomness
\end{itemize}

\medskip

\textbf{In contrast:}

If \(\mathbb{P}(A_n)=\frac{1}{n^2}\), then
\[
\sum_{n=1}^\infty \mathbb{P}(A_n)<\infty
\quad\Longrightarrow\quad
A_n \text{ occurs only finitely many times a.s.}
\]

\end{frame}

\begin{frame}{Borel--Cantelli Lemma}

Let \((A_n)_{n\ge1}\) be a sequence of events.
\pause
\begin{definitionblock}{Borel--Cantelli Lemma (II)}
Assume that the events \((A_n)\) are \textbf{independent} and that
$\sum_{n=1}^\infty \mathbb{P}(A_n) = \infty.$
\pause
Then
\[
\mathbb{P}\big(A_n \text{ infinitely often}\big) = 1.
\]
\end{definitionblock}

\medskip

\textbf{Interpretation:}

If independent events are frequent enough  
(non-summable probabilities),  
then infinitely many of them occur almost surely.

\end{frame}

\begin{frame}{Example: Borel--Cantelli Lemma (II)}

\textbf{Example — Coin tossing}

Consider an infinite sequence of independent coin tosses.

Define the events:
\[
A_n = \{\text{the $n$-th toss is Head}\}.
\]
\pause
Then:
\[
\mathbb{P}(A_n) = \frac{1}{2},
\qquad
\sum_{n=1}^\infty \mathbb{P}(A_n) = \infty.
\]

\medskip

Since the events \((A_n)\) are independent,  
Borel--Cantelli (II) applies and yields:
\[
\mathbb{P}(A_n \text{ infinitely often}) = 1.
\]

\medskip

\textbf{Conclusion:}
with probability one, Heads occur infinitely many times.

\end{frame}


% ======================================
% SECTION 2 — Random Variables & Distributions
% ======================================
\section{Random Variables \& Distributions}

\begin{frame}{Outline}
  \tableofcontents[currentsection, hideothersubsections]
\end{frame}

\subsection{Random Variables \& Distribution Functions}

\begin{frame}{Random Variables}
\begin{definitionblock}{Random Variable}
On a probability space \((\Omega,\mathbb{P})\), a (real-valued) random variable is a function
\[
  X:\Omega \to \mathbb{R}.
\]
\pause
\end{definitionblock}
Examples:
\begin{itemize}
  \item dice roll: \(\Omega=\{1,\dots,6\}\), \(X(\omega)=\omega\).\pause
  \item Indicator of an event \(A\): 
    \[
      \mathbf{1}_A(\omega) =
      \begin{cases}
        1 & \text{if } \omega\in A,\\
        0 & \text{otherwise.}
      \end{cases}
    \]\pause
\end{itemize}
\end{frame}

\begin{frame}{Distribution \& CDF}
\begin{definitionblock}{Distribution of \(X\)}
The \emph{distribution} (or \emph{law}) of \(X\), denoted \(\mathbb{P}_X\), is the probability measure on \(\mathbb{R}\) defined by
\[
  \mathbb{P}_X(A) = \mathbb{P}(X\in A)
  = \mathbb{P}(\{\omega : X(\omega)\in A\}).
\]
We also write \(X \sim \mathbb{P}_X\).\pause
\end{definitionblock}
\begin{definitionblock}{Cumulative Distribution Function}
The cdf of \(X\) is
\[
  F_X(t) = \mathbb{P}(X\le t),\quad t\in\mathbb{R}.
\]
\end{definitionblock}
Two random variables have the same distribution if and only if they have the same cdf.\pause
\end{frame}

\begin{frame}{Properties of the CDF}
For any random variable \(X\):\pause
\begin{itemize}
  \item \(F_X\) is non-decreasing;\pause
  \item \(F_X\) is right-continuous;\pause
  \item \(\displaystyle \lim_{t\to -\infty} F_X(t)=0\), \(\displaystyle \lim_{t\to +\infty} F_X(t)=1\).\pause
\end{itemize}
Conversely, any function \(F:\mathbb{R}\to[0,1]\) with these properties
is a cdf of some random variable.
\end{frame}

\begin{frame}{Example of a CDF (Discrete Random Variable)}

Let \(X\) be the outcome of a fair dice:
$
\mathbb{P}(X=k)=\frac{1}{6},\quad k=1,\dots,6.
$
The cumulative distribution function is
$
F_X(t)=\mathbb{P}(X\le t).
$\pause
Explicitly:
\[
F_X(t)=
\begin{cases}
0, & t<1,\\
\frac{1}{6}, & 1\le t<2,\\
\frac{2}{6}, & 2\le t<3,\\
\frac{3}{6}, & 3\le t<4,\\
\frac{4}{6}, & 4\le t<5,\\
\frac{5}{6}, & 5\le t<6,\\
1, & t\ge6.
\end{cases}
\vspace{-9mm}
\]

\textbf{Observation:}
\begin{itemize}
  \item step function
  \item non-decreasing
  \item right-continuous
\end{itemize}

\end{frame}

\begin{frame}{Example of a CDF (Continuous Random Variable)}

Let \(X\sim U[0,1]\).

The density is
\[
f(x)=\mathbf{1}_{[0,1]}(x),
\]\pause
and the CDF is
\[
F_X(t)=
\begin{cases}
0, & t<0,\\
t, & 0\le t\le1,\\
1, & t>1.
\end{cases}
\]\pause

\medskip

\textbf{Observation:}
\begin{itemize}
  \item continuous
  \item increasing
  \item limits: \(0\) at \(-\infty\), \(1\) at \(+\infty\)
\end{itemize}

\end{frame}


\subsection{Discrete \& Continuous Distributions}

\begin{frame}{Discrete Random Variables}
\begin{definitionblock}{Discrete Distribution}
\(X\) is \emph{discrete} if it takes values in a finite or countable set
\(\{x_1,x_2,\dots\}\).\pause
\end{definitionblock}
\begin{itemize}
  \item Probability mass function (pmf):
    \[
      p(x_k) = \mathbb{P}(X=x_k),\quad x_k\in\{x_1,x_2,\dots\}.
    \]\pause
  \item Cdf:
    \[
      F_X(t) = \sum_{x_k \le t} p(x_k).
    \]
\end{itemize}
\end{frame}

\begin{frame}{Examples of Discrete Laws}
\begin{itemize}
  \item \textbf{Bernoulli} \(B(p)\), parameter \(p\in[0,1]\), values \(\{0,1\}\):
  \[
    \mathbb{P}(X=1)=p,\quad \mathbb{P}(X=0)=1-p.
  \]\pause
  \item \textbf{Binomial} \(B(n,p)\), parameters \(n\in\mathbb{N}, p\in[0,1]\):
  \[
    \mathbb{P}(X = k) = \binom{n}{k} p^k (1-p)^{n-k},
    \quad k=0,\dots,n.
  \]\pause
  \item \textbf{Geometric}, parameter \(p\in(0,1]\):
  \[
    \mathbb{P}(X=k) = (1-p)^{k-1} p,\quad k=1,2,\dots
  \]\pause
  \item \textbf{Poisson}, parameter \(\lambda>0\):
  \[
    \mathbb{P}(X=k) = e^{-\lambda}\frac{\lambda^k}{k!},\quad k=0,1,2,\dots
  \]\pause
\end{itemize}
\end{frame}

\begin{frame}{Shape of the CDF for Discrete Laws}
For a discrete random variable:\pause
\begin{itemize}
  \item The cdf \(F_X\) is a \textbf{step function}.\pause
  \item Jumps occur at the values actually taken by \(X\).\pause
  \item The height of the jump at \(x_k\) equals \(\mathbb{P}(X=x_k)\).\pause
\end{itemize}

\medskip

\textbf{Example:} for a Bernoulli\((p)\),
\[
X\in\{0,1\},\quad
F_X(t)=
\begin{cases}
0 & t<0,\\
1-p & 0\le t<1,\\
1 & t\ge 1.
\end{cases}
\]

\end{frame}

\begin{frame}{Continuous Random Variables}
\begin{definitionblock}{Density}
\(X\) has a \emph{continuous} distribution if there exists a non-negative
function \(f:\mathbb{R}\to[0,\infty)\) such that for all Borel sets \(A\),
\[
  \mathbb{P}(X\in A) = \int_A f(x)\,dx,
\]
with \(\displaystyle\int_{\mathbb{R}} f(x)\,dx = 1\).\pause
The function \(f\) is called the \emph{density} of \(X\).
\end{definitionblock}
Then
\[
  F_X(t) = \int_{-\infty}^t f(x)\,dx, \qquad
  f(t) = F_X'(t)\ \text{for almost all }t.
\]
\end{frame}

\begin{frame}{Examples of Continuous Laws}
\begin{itemize}
  \item \textbf{Uniform} \(U[a,b]\):
  \[
    f(x) = \frac{1}{b-a}\mathbf{1}_{[a,b]}(x).
  \]\pause
  \item \textbf{Exponential} \(\mathrm{Exp}(\lambda)\), \(\lambda>0\):
  \[
    f(x) = \lambda e^{-\lambda x}\mathbf{1}_{x\ge 0}.
  \]\pause
  \item \textbf{Normal} (Gaussian) \(\mathcal{N}(\mu,\sigma^2)\), \(\sigma^2>0\):
  \[
    f(x) = \frac{1}{\sigma\sqrt{2\pi}}
    \exp\left(-\frac{(x-\mu)^2}{2\sigma^2}\right).
  \]\pause
\end{itemize}
\end{frame}

\begin{frame}{CDF of Continuous Distributions}
For a continuous random variable with density \(f\):
\begin{itemize}
  \item The cdf is
    \[
      F_X(t)=\int_{-\infty}^t f(x)\,dx.
    \]\pause
  \item \(F_X\) is \textbf{continuous} on \(\mathbb{R}\).\pause
  \item At almost every point \(t\), \(f(t)=F_X'(t)\).\pause
\end{itemize}

\medskip

\textbf{Examples:}
\begin{itemize}
  \item Exponential distribution: increasing, concave cdf on \([0,\infty)\).
  \item Normal distribution: smooth, S-shaped cdf on \(\mathbb{R}\).
\end{itemize}
\end{frame}

\begin{frame}{Typical Shapes of Continuous CDFs}

\begin{center}
\begin{tikzpicture}[scale=0.9]

% Axes for Exponential
\draw[->] (-0.2,0) -- (4.2,0) node[right] {$t$};
\draw[->] (0,-0.2) -- (0,3) node[above] {$F_X(t)$};
\node at (2,-0.6) {\small Exponential CDF};

% Exponential CDF curve
\draw[thick, BootcampBlueDark, domain=0:4, smooth]
  plot (\x,{2.5*(1-exp(-0.8*\x))});

% Axes for Normal
\begin{scope}[xshift=6cm]
\draw[->] (-3.2,0) -- (3.2,0) node[right] {$t$};
\draw[->] (0,-0.2) -- (0,3) node[above] {$F_X(t)$};
\node at (0,-0.6) {\small Normal CDF};

% Normal CDF (stylized S-shape)
\draw[thick, AccentGold, smooth]
  plot coordinates {
    (-3,0.1) (-2,0.3) (-1,0.8) (0,1.5)
    (1,2.2) (2,2.7) (3,2.9)
  };
\end{scope}

\end{tikzpicture}
\end{center}

\medskip
\pause
\textbf{Key visual features:}
\begin{itemize}
  \item CDFs are increasing and bounded between 0 and 1
  \item Exponential: fast increase, then saturation
  \item Normal: smooth S-shape, symmetric around the mean
\end{itemize}

\end{frame}

\begin{frame}{Beyond Discrete and Continuous}
Not all distributions are purely discrete or purely continuous.

\medskip
\pause
\textbf{Example (censored exponential).}
Let \(Y \sim \mathrm{Exp}(1)\) and define
\[
X = \min\{1,Y\}.
\]\pause
Then:
\begin{itemize}
  \item On \([0,1)\), \(X\) has a density inherited from \(Y\).
  \item At the point \(1\), there is an \textbf{atom}:
    \(\mathbb{P}(X=1) > 0\).
\end{itemize}

\medskip
\pause
\textbf{Conclusion:} \(X\) has a \emph{mixed} distribution
(neither purely discrete nor purely continuous).

\end{frame}

\begin{frame}{CDF of a Mixed Distribution: Censored Exponential}

Let \(Y\sim \mathrm{Exp}(1)\) and \(X=\min\{1,Y\}\).
The cdf of \(X\) is
\[
F_X(t)=
\begin{cases}
0, & t<0,\\[1mm]
1-e^{-t}, & 0\le t<1,\\[1mm]
1, & t\ge 1,
\end{cases}
\vspace{-5mm}
\]
so there is a jump at \(t=1\):
$
\mathbb{P}(X=1)=1-F_X(1^-)=e^{-1}.
$
\begin{center}
\begin{tikzpicture}[scale=0.9]

% =========================
% Left panel: continuous part on [0,1)
% =========================
\begin{scope}
  % axes (left)
  \draw[->] (-0.2,0) -- (2.4,0) node[right] {$t$};
  \draw[->] (0,-0.2) -- (0,3.2) node[above] {$F_X(t)$};

  % curve on [0,1)
  \draw[thick, BootcampBlueDark, domain=0:2, smooth]
    plot (\x,{2.6*(1-exp(-\x/2))});

  % mark t=2 corresponds to t=1 in original scale
  \draw[dashed] (2,0) -- (2,3.1);
  \node[below] at (2,0) {\small $1^{-}$};

  % open circle at the right endpoint (left limit)
  \draw[BootcampBlueDark, thick] (2,{2.6*(1-exp(-1/2))}) circle (3pt);
  \filldraw[white] (2,{2.6*(1-exp(-1/2))}) circle (2pt);

  \node at (1.2,2.9) {\small on $[0,1)$};
\end{scope}

% =========================
% Right panel: jump at 1 and constant after
% =========================
\begin{scope}[xshift=6cm]
  % axes (right)
  \draw[->] (-0.2,0) -- (2.6,0) node[right] {$t$};
  \draw[->] (0,-0.2) -- (0,3.2) node[above] {$F_X(t)$};

  % show the jump at t=1
  \draw[dashed] (0.9,0) -- (0.9,3.1);
  \node[below] at (0.9,0) {\small $1$};

  % left-limit value (open circle)
  \draw[BootcampBlueDark, thick] (0.9,1.9) circle (3pt);
  \filldraw[white] (0.9,1.9) circle (2pt);

  % jump up to 1
  \draw[thick, BootcampBlueDark] (0.9,1.9) -- (0.9,3);

  % value at t=1 (closed dot)
  \filldraw[BootcampBlueDark] (0.9,3) circle (2.5pt);

  % constant part for t>=1
  \draw[thick, BootcampBlueDark] (0.95,3) -- (2.4,3);

  \node[align=center] at (1.6,2.4) {\small jump \& constant\\[-1mm]\small for $t\ge 1$};
\end{scope}

% Optional: a light separator between panels
\draw[gray!50] (3.2,-0.2) -- (3.2,3.2);

\end{tikzpicture}
\end{center}

\end{frame}

\subsection{Expectation, Variance, Moments}

\begin{frame}{Expectation}
\begin{definitionblock}{Expectation}
Let \(X\) be a random variable.
\begin{itemize}
  \item If \(X\) is discrete with values \(\{x_k\}\),
  \[
    \mathbb{E}[X] = \sum_{k\ge1} x_k\,\mathbb{P}(X=x_k),
  \]\pause
  provided \(\sum_{k\ge1} |x_k|\,\mathbb{P}(X=x_k) < \infty\).
  \item If \(X\) is continuous with density \(f\),
  \[
    \mathbb{E}[X] = \int_{-\infty}^{\infty} x\,f(x)\,dx,
  \]\pause
  provided \(\int_{-\infty}^{\infty} |x|f(x)\,dx < \infty\).
\end{itemize}
\end{definitionblock}
\end{frame}

\begin{frame}{Example: Expectation (Discrete Case)}

\textbf{Example 1 — Fair dice}

Let \(X\) be the outcome of a fair dice:
\[
\mathbb{P}(X=k)=\frac{1}{6}, \quad k=1,\dots,6.
\]
\pause
Then
\[
\mathbb{E}[X]
= \sum_{k=1}^6 k \cdot \frac{1}{6}
= \frac{1+2+3+4+5+6}{6}
= \frac{21}{6}
= 3.5.
\]\pause

\medskip

\textbf{Example 2 — Bernoulli random variable}

If \(X\sim B(p)\), then
\[
\mathbb{E}[X]
= 1\cdot p + 0\cdot(1-p)
= p.
\]\pause

\end{frame}

\begin{frame}{Example: Expectation (Continuous Case)}

\textbf{Example — Exponential distribution}

Let \(X \sim \mathrm{Exp}(\lambda)\), with density
\[
f(x)=\lambda e^{-\lambda x}\mathbf{1}_{x\ge0}.
\]
\pause
Then
\[
\mathbb{E}[X]
= \int_0^\infty x\,\lambda e^{-\lambda x}\,dx
= \frac{1}{\lambda}.
\]
\pause
\medskip

\textbf{Example — Uniform distribution}

If \(X\sim U[a,b]\), then
\[
\mathbb{E}[X]
= \int_a^b x\,\frac{1}{b-a}\,dx
= \frac{a+b}{2}.
\]

\end{frame}

\begin{frame}{When the Expectation Does Not Exist}
The expectation is not defined for all distributions.

\medskip

\textbf{Example: Cauchy distribution.}
\[
f(x)=\frac{1}{\pi(1+x^2)},\quad x\in\mathbb{R}.
\]
Then
\[
\int_{-\infty}^{\infty} |x| f(x)\,dx = \infty,
\]
so \(\mathbb{E}[X]\) does \emph{not} exist.

\medskip

\textbf{Lesson:} finite expectation requires \(\mathbb{E}[|X|]<\infty\).

\end{frame}

\begin{frame}{Properties of Expectation}
For random variables \(X,Y\) and constants \(a,b\):\pause
\begin{itemize}
  \item Linearity:
    \[
      \mathbb{E}[aX + bY] = a\,\mathbb{E}[X] + b\,\mathbb{E}[Y].
    \]\pause
  \item Monotonicity: if \(X \le Y\) almost surely, then
    \(\mathbb{E}[X] \le \mathbb{E}[Y]\).\pause
  \item If \(X = \mathbf{1}_A\), then \(\mathbb{E}[X] = \mathbb{P}(A)\).\pause
\end{itemize}
\end{frame}

\begin{frame}{Functions of a Random Variable}
Let \(Y = \varphi(X)\) for some measurable \(\varphi:\mathbb{R}\to\mathbb{R}\).
If \(\mathbb{E}[|Y|]<\infty\), then:\pause
\begin{itemize}
  \item Discrete case:
  \[
    \mathbb{E}[\varphi(X)] = \sum_k \varphi(x_k)\,\mathbb{P}(X=x_k).
  \]\pause
  \item Continuous case:
  \[
    \mathbb{E}[\varphi(X)] = \int_{\mathbb{R}} \varphi(x)\,f(x)\,dx.
  \]\pause
\end{itemize}
\end{frame}

\begin{frame}{Example: Functions of a Random Variable}

\textbf{Discrete example}

Let \(X\sim B(p)\) (Bernoulli random variable) and define
$
Y=\varphi(X)=X^2.
$
\pause
Since \(X\in\{0,1\}\), we have \(X^2=X\), hence
\[
\mathbb{E}[X^2]
= \sum_{x\in\{0,1\}} x^2\,\mathbb{P}(X=x)
= 0^2(1-p)+1^2p
= p.
\]
\pause
\medskip

\textbf{Continuous example}

Let \(X\sim U[0,1]\) and define
\[
Y=\varphi(X)=X^2.
\]
\pause
Then
\[
\mathbb{E}[X^2]
= \int_0^1 x^2\,dx
= \frac{1}{3}.
\]
\pause
\medskip

\textbf{Key idea:}  
we do not need the distribution of \(Y\); we integrate \(\varphi(X)\) directly.

\end{frame}

\begin{frame}{Non-negative Variables with Zero Mean}
\begin{definitionblock}{A Useful Property}
Let \(X\) be a non-negative random variable:
\[
\mathbb{P}(X\ge 0)=1.
\]
If \(\mathbb{E}[X]=0\), then
\[
X=0 \quad \text{almost surely}.
\]
\end{definitionblock}
\pause
\medskip

\textbf{Intuition:}  
a non-negative variable can have zero average only if it is zero
with probability \(1\).

\end{frame}

\begin{frame}{Variance and Higher Moments}
\begin{definitionblock}{Variance}
If \(\mathbb{E}[X^2]<\infty\), the variance of \(X\) is
\[
  \mathrm{Var}(X) = \mathbb{E}\big[(X-\mathbb{E}[X])^2\big]
  = \mathbb{E}[X^2] - (\mathbb{E}[X])^2.
\]\pause
The standard deviation is \(\sigma_X = \sqrt{\mathrm{Var}(X)}\).
\end{definitionblock}
\begin{itemize}
  \item \(\mathrm{Var}(X)\ge0\), and \(\mathrm{Var}(X)=0\) iff \(X=c\) a.s. for some constant \(c\).\pause
  \item For constants \(a,b\): \(\mathrm{Var}(aX+b)=a^2\mathrm{Var}(X)\).\pause
\end{itemize}
\medskip
\textbf{\(k\)-th moment:} if \(\mathbb{E}[|X|^k]<\infty\), then \(\mathbb{E}[X^k]\).  \pause
\textbf{\(k\)-th central moment:} \(\mathbb{E}[(X-\mathbb{E}[X])^k]\).\pause
\end{frame}


\subsection{Inequalities \& Tools}

\begin{frame}{Useful Inequalities}
\begin{itemize}
  \item \textbf{Jensen's inequality.}  
  Let \(\varphi:\mathbb{R}\to\mathbb{R}\) be convex and \(X\) integrable:\pause
  \[
    \varphi(\mathbb{E}[X]) \le \mathbb{E}[\varphi(X)].
  \]\pause
  \item \textbf{Markov's inequality.}  
  If \(X\ge 0\) a.s. and \(c>0\),\pause
  \[
    \mathbb{P}(X \ge c) \le \frac{\mathbb{E}[X]}{c}.
  \]\pause
  \item \textbf{Chebyshev's inequality.}  
  If \(\mathbb{E}[X^2]<\infty\) and \(a>0\),\pause
  \[
    \mathbb{P}(|X-\mathbb{E}[X]|\ge a)
    \le \frac{\mathrm{Var}(X)}{a^2}.
  \]\pause
\end{itemize}
\end{frame}

\begin{frame}{Examples of Useful Inequalities}

\textbf{Jensen's inequality}

Let \(X\sim U[0,1]\) and \(\varphi(x)=x^2\) (convex function).
\[
\varphi(\mathbb{E}[X]) = \left(\tfrac12\right)^2 = \tfrac14,
\qquad
\mathbb{E}[\varphi(X)] = \mathbb{E}[X^2] = \tfrac13.
\]
Thus $
\varphi(\mathbb{E}[X]) \le \mathbb{E}[\varphi(X)].
$

\medskip

\textbf{Markov's inequality}

Let \(X\ge0\) with \(\mathbb{E}[X]=2\). For \(c=4\),
\[
\mathbb{P}(X\ge4) \le \frac{2}{4}=\tfrac12.
\]
This gives a rough upper bound using only the mean.

\medskip

\textbf{Chebyshev's inequality}

Let \(X\) with \(\mathbb{E}[X]=0\) and \(\mathrm{Var}(X)=1\).
Then for \(a=2\),
\[
\mathbb{P}(|X|\ge2) \le \frac{1}{4}.
\]

\medskip

\textbf{Key message:}
these inequalities give bounds without knowing the full distribution.

\end{frame}

\begin{frame}{How to Find the Distribution of \(X\)?}
Some common methods:\pause
\begin{itemize}
  \item Compute its cdf \(F_X(t) = \mathbb{P}(X\le t)\) directly, then differentiate to get the density (if it exists).\pause
  \item Use a change of variable in integrals to find the density of a transformation, e.g. \(Y=aX+b\) or \(Y=g(X)\).\pause
  \item Use the \emph{characteristic function}
  \[
    \Phi_X(t) = \mathbb{E}[e^{itX}],\quad t\in\mathbb{R}.
  \]\pause
  If \(\Phi_X = \Phi_Y\), then \(X\) and \(Y\) have the same law.
\end{itemize}
\end{frame}

\begin{frame}{Example: Finding a Distribution via the CDF}

Let \(X\sim U[0,1]\) and define
$
Y = X^2.
$
\pause

\textbf{Step 1: Compute the CDF of \(Y\).}
For \(t\in[0,1]\),
$
F_Y(t)
= \mathbb{P}(Y \le t)
= \mathbb{P}(X^2 \le t)
= \mathbb{P}(X \le \sqrt{t})
= \sqrt{t}.
$
\pause
Moreover,
\[
F_Y(t)=
\begin{cases}
0, & t<0,\\
\sqrt{t}, & 0\le t\le 1,\\
1, & t\ge 1.
\end{cases}
\vspace{-5mm}
\]
\pause
\medskip

\textbf{Step 2: Differentiate to get the density.}

For \(t\in(0,1)\),
$
f_Y(t) = F_Y'(t) = \frac{1}{2\sqrt{t}}.
$
\pause

\textbf{Conclusion:}
\(Y\) has a continuous distribution on \([0,1]\) with density
\[
f_Y(t)=\frac{1}{2\sqrt{t}}\,\mathbf{1}_{(0,1)}(t).
\]
\pause
\end{frame}

\subsection{\(L^p\) Spaces and Convergence}

\begin{frame}{\(L^p\) Spaces}
\begin{definitionblock}{\(L^p\) Space}
For \(p\ge1\), \(L^p(\Omega,\mathbb{P})\) is the set of random variables \(X\) such that \(\mathbb{E}[|X|^p]<\infty\). The \(L^p\) norm is
\[
  \|X\|_p = \big(\mathbb{E}[|X|^p]\big)^{1/p}.
\]\pause
\end{definitionblock}
\begin{itemize}
  \item \(L^p\) is a vector space.
  \item \(\|\cdot\|_p\) is a norm (triangle inequality, etc.).
  \item For \(q>p\), we have \(L^q \subset L^p\).
\end{itemize}
\end{frame}

\begin{frame}{The Space \(L^2\) and Cauchy--Schwarz}
\begin{itemize}
  \item If \(X,Y\in L^2\), then the inner product
    \[
      \langle X,Y\rangle = \mathbb{E}[XY]
    \]
    makes \(L^2\) a Hilbert space.
  \item \textbf{Cauchy--Schwarz inequality:}
    \[
      |\mathbb{E}[XY]| \le
      \big(\mathbb{E}[X^2]\big)^{1/2}
      \big(\mathbb{E}[Y^2]\big)^{1/2}.
    \]\pause
  \item Orthogonality: \(X\perp Y\) if \(\langle X,Y\rangle=0\).\pause
\end{itemize}
\end{frame}

\begin{frame}{Swapping Expectation and Limit (Preview)}
Let \((X_n)_{n\ge1}\) be a sequence of random variables and \(X\) a limit.\pause
\begin{itemize}
  \item \textbf{Monotone convergence theorem.}  
  If \(X_n\ge0\) and \(X_n\uparrow X\) almost surely, then
  \[
    \mathbb{E}[X] = \lim_{n\to\infty} \mathbb{E}[X_n].
  \]\pause
  \item \textbf{Dominated convergence theorem.}  
  If \(X_n\to X\) a.s. and \(|X_n|\le Y\) for some integrable \(Y\), then
  \[
    \mathbb{E}[X] = \lim_{n\to\infty} \mathbb{E}[X_n].
  \]\pause
\end{itemize}
These results justify interchanging limits and expectations in many situations.
\end{frame}

\begin{frame}{Application: Swapping Sum and Expectation}
Let \((X_k)_{k\ge0}\) be a sequence of non-negative random variables.
\begin{itemize}
  \item Define the partial sums
  \[
  S_n = \sum_{k=0}^n X_k,\quad n\ge0.
  \]\pause
  \item Then \(S_n\uparrow S=\sum_{k=0}^\infty X_k\) (possibly \(+\infty\)).
\end{itemize}

By the monotone convergence theorem:
\[
\mathbb{E}\Big[\sum_{k=0}^\infty X_k\Big]
= \sum_{k=0}^\infty \mathbb{E}[X_k],
\]\pause
where both sides may be infinite.

\medskip

\textbf{Use case:} series of payoffs, risk contributions, etc.

\end{frame}
% ======================================
% SECTION 3 — Summary
% ======================================

\end{document}